\documentclass[11pt]{article}
\usepackage{amsmath,amsthm,amssymb}
\usepackage[margin=1in]{geometry}
\usepackage{hyperref}
\usepackage{enumitem}
\usepackage{booktabs}

\newtheorem{theorem}{Theorem}
\newtheorem{lemma}[theorem]{Lemma}
\newtheorem{proposition}[theorem]{Proposition}
\newtheorem{corollary}[theorem]{Corollary}
\newtheorem{conjecture}[theorem]{Conjecture}
\theoremstyle{definition}
\newtheorem{definition}[theorem]{Definition}
\newtheorem{remark}[theorem]{Remark}
\newtheorem{example}[theorem]{Example}

\title{The ``factor-of-two'' conjecture for $\dim B^{(\lambda)}_{\ell+1} V_\lambda$ \\
  is too coarse: shape-dependent non-stability}
\author{Clio}
\date{14 May 2026 (evening prove session)}

\begin{document}
\maketitle

\begin{abstract}
The May-13/14 PROVE.md note proposed a side conjecture for the
``factor-of-two sub-family'': partitions $\lambda$ with two equal
non-trivial rows ($\lambda_i = \lambda_{i+1} \ge 2$ for some $i$)
should satisfy $\dim B_{\ell+1}^{(\lambda)} V_\lambda = 2 \cdot
f^{\lambda^{(\ge 2)}}$, with two cited examples: $(4,2,2,1)$ giving
$12 = 2{\cdot}6$ and $(4,4,1,1)$ giving $10 = 2{\cdot}5$.

We verify the two cited examples and \textbf{refute the
generalisation} via systematic computation:
\begin{itemize}[topsep=0.3em,itemsep=0.2em]
\item $(6, 2, 1)$ at $n=9$ has $\dim B = 17$, neither stable ($5$) nor
  $2{\cdot}5$.  No row-repeat is even present here.
\item $(3, 3, 1)$ at $n=7$ has $\dim B = 3 = 1.5{\cdot}f^{(2,2)}$ — a
  row-repeat at the top but factor $\ne 2$.
\item $(4, 3, 1, 1)$ at $n=9$ has $\dim B = 12 = 2.4{\cdot}f^{(3,2)}$ —
  no row-repeat, factor not in $\{1, 2\}$.
\item $(4, 4, 1, 1, 1)$ at $n=11$ has $\dim B = 15 = 3{\cdot}5$ — same
  row-repeat as $(4,4,1,1)$ but factor $3$, not $2$.
\end{itemize}

The actual picture is shape- and $q$-dependent: each shape
$\widehat\lambda$ (non-trivial part) has a stability threshold
$q_0(\widehat\lambda)$ such that $\lambda = (\widehat\lambda, 1^q)$
satisfies the May-13 closed form $\dim B = f^{\lambda^{(\ge 2)}}$ for
$q \ge q_0$, with no clean closed form for $q < q_0$.  Tabulated data
(\S\ref{sec:data}) shows the non-stable values range from
sub-stable ($\dim B < f^{\lambda^{(\ge 2)}}$) through strongly
super-stable (factor $> 3$).

The refined conjecture is therefore: \emph{$f^{\lambda^{(\ge 2)}}$
is the closed form precisely above a shape-dependent threshold; below
the threshold, the dimension is genuinely a finer invariant of
$\lambda$.}  We close with the threshold data and an open question
about its closed form.
\end{abstract}

\section{Setup}

We use the conventions of the May-13--May-20 series.  Let
$H_q(S_n)$ be the Iwahori--Hecke algebra over $\mathbb{Q}(q)$ with
generators $T_1, \ldots, T_{n-1}$, quadratic relation $T_i^2 = (q-1)T_i
+ q$, and braid relations.  For $\lambda \vdash n$ write $V_\lambda$ for
the Specht module in the Hoefsmit seminormal basis indexed by
$\mathrm{SYT}(\lambda)$.  Write
\[
   R'_j := (T_{j-1}{+}1)(T_{j-2}{+}1)\cdots(T_1{+}1)
   \quad\in\, H_q(S_j),
   \qquad
   B_j^{(\lambda)} V_\lambda := R'_j R'_{j+1}\cdots R'_n V_\lambda.
\]
The \emph{sharp index} is $j_0 = \ell(\lambda) + 1$, and the operator
of interest is $B_{\ell+1}^{(\lambda)} V_\lambda$.

For $\lambda$ with $\lambda_\ell = 1$, write $\lambda =
(\widehat\lambda, 1^q)$ where $\widehat\lambda$ is the
``non-trivial part'' (rows of length $\ge 2$) and $q = q(\lambda)$ is
the number of \emph{trailing $1$-rows}.  The \emph{decoration shape} is
$\lambda^{(\ge 2)} := (\lambda_1 - 1, \lambda_2 - 1, \ldots)$, obtained
from $\lambda$ by deleting column~$1$.

\begin{theorem}[May-13 stable closed form, conjectural]
\label{thm:stable}
There exists a shape-dependent threshold $q_0(\widehat\lambda) \ge 0$
such that for $q \ge q_0(\widehat\lambda)$,
\[
   \dim B_{\ell + 1}^{(\lambda)} V_\lambda \;=\; f^{\lambda^{(\ge 2)}}.
\]
This has been verified on $9$ families with $q$ exceeding small
threshold values.
\end{theorem}

\section{The PROVE.md conjecture}\label{sec:prove-md}

The May-13/14 PROVE.md side-question conjectured the following.

\begin{conjecture}[Factor-of-two, PROVE.md side question]\label{conj:factor2}
For partitions $\lambda$ with at least one pair of equal non-trivial
rows ($\exists\, i : \lambda_i = \lambda_{i+1} \ge 2$),
\[
   \dim B_{\ell + 1}^{(\lambda)} V_\lambda \;=\; 2 \cdot f^{\lambda^{(\ge 2)}}.
\]
Cited supporting examples:
\begin{itemize}
\item $(4, 2, 2, 1)$ at $n = 9$: $\dim B = 12 = 2 \cdot f^{(3,1,1)} = 2{\cdot}6$.
\item $(4, 4, 1, 1)$ at $n = 10$: $\dim B = 10 = 2 \cdot f^{(3,3)} = 2{\cdot}5$.
\end{itemize}
\end{conjecture}

\paragraph{Both cited examples are computationally confirmed.}
At $Q = 11$ mod $P = 100\,003$ via the Hoefsmit seminormal basis (script
\texttt{\textasciitilde/projects/scratch/2026-05-14-row-repeat/test\_base.py}):
\[
   \dim B^{(4,2,2,1)}_5 V_{(4,2,2,1)} \;=\; 12, \qquad
   \dim B^{(4,4,1,1)}_5 V_{(4,4,1,1)} \;=\; 10.
\]

\section{Refutation: four counter-examples}\label{sec:refutation}

Conjecture~\ref{conj:factor2} as stated is false.  We exhibit four
explicit counter-examples covering distinct failure modes.

\subsection{$(6, 2, 1)$ at $n = 9$: factor $3.4$, no row-repeat}

The partition $\lambda = (6, 2, 1)$ has \emph{no} row-repeat
($\lambda_1 = 6 \ne 2 = \lambda_2$), so it lies outside the
conjecture's scope.  Yet it illustrates that even \emph{without} a
row-repeat, $\dim B$ at $q = 1$ can deviate from both the stable form
and the proposed factor-of-two:
\[
   \dim V_{(6,2,1)} = 105, \qquad f^{(5,1)} = 5, \qquad
   \dim B^{(6,2,1)}_4 V_{(6,2,1)} = 17.
\]
Here $17 / 5 = 3.4$ is non-integer.  Confirmed at $Q \in \{11, 13, 17,
23\}$, so this is not an artefact of the mod-$P$ reduction.  This
proves that ``factor of two'' is not even the correct \emph{integer}
factor in the family $(k, 2, 1)$ at $k \ge 6$.

\subsection{$(3, 3, 1)$ at $n = 7$: row-repeat at the top, factor $1.5$}

$\lambda = (3, 3, 1)$ has a row-repeat ($\lambda_1 = \lambda_2 = 3 \ge
2$).  Its decoration shape is $(2, 2)$, $f^{(2,2)} = 2$.  The
conjecture predicts $\dim B = 4$.  Actual:
\[
   \dim B^{(3,3,1)}_4 V_{(3,3,1)} = 3.
\]
$3/2 = 1.5$; in particular, $3 \ne 4$.  This refutes the conjecture for
the simplest row-repeat shape with $q = 1$.

\subsection{$(4, 3, 1, 1)$ at $n = 9$: factor $2.4$}

$\lambda = (4, 3, 1, 1)$ has no row-repeat.  Decoration $(3, 2)$,
$f^{(3, 2)} = 5$.  Actual:
\[
   \dim B^{(4,3,1,1)}_5 V_{(4,3,1,1)} = 12 = 2.4 \cdot 5.
\]
A non-integer factor between $2$ and $3$.

\subsection{$(4, 4, 1, 1, 1)$ at $n = 11$: same row-repeat, factor $3$}

$\lambda = (4, 4, 1, 1, 1)$ has the same row-repeat as the
PROVE.md example $(4, 4, 1, 1)$ but one more trailing $1$-row.
Decoration $(3, 3)$, $f^{(3, 3)} = 5$.  Actual:
\[
   \dim B^{(4,4,1,1,1)}_6 V_{(4,4,1,1,1)} = 15 = 3 \cdot 5.
\]
So the conjecture's ``factor of two'' is actually a factor of three at
$q = 3$.  By scanning $(4, 4, 1^q)$ for varying $q$ (data below),
the factor descends as $q$ increases beyond a critical $q_0((4, 4))$.

\section{The actual picture: shape-dependent stability threshold}
\label{sec:picture}

The right framework, supported by the data of \S\ref{sec:data}, is:

\begin{conjecture}[Refined stability]\label{conj:refined}
For each partition $\widehat\lambda$ with $\widehat\lambda_\ell \ge 2$,
there exists a threshold $q_0(\widehat\lambda) \ge 0$ such that for
$\lambda = (\widehat\lambda, 1^q)$:
\begin{itemize}[topsep=0.3em,itemsep=0.2em]
\item For $q \ge q_0(\widehat\lambda)$ (the \emph{stable regime}),
$\dim B_{\ell + 1}^{(\lambda)} V_\lambda = f^{\lambda^{(\ge 2)}}$.
\item For $q < q_0(\widehat\lambda)$ (the \emph{non-stable regime}),
$\dim B_{\ell + 1}^{(\lambda)} V_\lambda$ is a finer invariant of
$\widehat\lambda$ and $q$, with no closed form in terms of
$f^{\lambda^{(\ge 2)}}$ alone.
\end{itemize}
The threshold $q_0(\widehat\lambda)$ depends on $\widehat\lambda$ in a
non-trivial way: the data of \S\ref{sec:data} place $q_0$ in the range
$\{0, 1, 2, 3, 4, \ldots\}$ depending on shape, with no obvious
closed form yet.
\end{conjecture}

\paragraph{What's known about $q_0$.}
\begin{itemize}[topsep=0.3em,itemsep=0.2em]
\item $q_0(\widehat\lambda) = 0$ for ``small'' shapes:
$(k)$ (hook); $(3, 2, 2)$, $(4, 2, 2)$; possibly others (all tested
$q = 0$ cases that are stable).
\item $q_0(\widehat\lambda) \ge 1$ for $\widehat\lambda \in \{(2, 2),
(3, 2), (4, 2), (3, 3), (3, 3, 2), \ldots\}$.
\item $q_0((4, 3)) \ge 3$ — the first family where the threshold is
$\ge 3$.
\item For large $\widehat\lambda_1$ (e.g.\ $(6, 2)$, $(5, 2)$), $q_0$
is at least $1$ and likely grows with $\widehat\lambda_1$ (data
incomplete past $n = 11$).
\end{itemize}

\paragraph{Why not a simple formula in $\widehat\lambda_1$.}
The candidate $q_0(\widehat\lambda) = \widehat\lambda_1 - 2$ would
give:
\begin{itemize}
\item $(2, 2) \mapsto 0$. Tested $q_0 \ge 1$. \textbf{Fail.}
\item $(3, 2) \mapsto 1$. Matches data.
\item $(4, 2) \mapsto 2$. Matches data.
\item $(3, 3) \mapsto 1$. Tested $q_0 = 2$. \textbf{Fail.}
\item $(4, 3) \mapsto 2$. Tested $q_0 = 3$. \textbf{Fail.}
\end{itemize}
The first failure shows $q_0$ depends on more than $\widehat\lambda_1$.
The $(3, 3)$ case suggests row-repeats \emph{do} affect $q_0$, just
not by the multiplicative factor that Conjecture~\ref{conj:factor2}
proposed.

\section{Data}\label{sec:data}

All $\dim B$ values are computed mod $P = 100\,003$ at $Q = 11$ via
\texttt{compute\_dim\_B.py} in
\texttt{\textasciitilde/projects/scratch/2026-05-13-multi-corner/}.

\subsection{The $(k, 2, 1)$ family at $q = 1$ (no row-repeat)}

\begin{center}
\begin{tabular}{c|cccc}
\toprule
$k$        & $3$ & $4$ & $5$ & $6$ \\
\midrule
$\lambda$  & $(3,2,1)$ & $(4,2,1)$ & $(5,2,1)$ & $(6,2,1)$ \\
$n$        & $6$       & $7$       & $8$       & $9$       \\
$f^{(k-1, 1)}$ & $2$   & $3$       & $4$       & $5$       \\
$\dim B$   & $2$       & $6$       & $8$       & $17$      \\
factor     & $1$       & $2$       & $2$       & $3.4$     \\
\bottomrule
\end{tabular}
\end{center}

This family shows three distinct regimes: stable at $k = 3$, factor
$2$ at $k = 4, 5$, factor $3.4$ at $k = 6$.

\subsection{The $(5, 2, 1^q)$ family across $q$ --- a clean linear pattern in non-stable regime}

A particularly clean phenomenon emerges in this family.

\begin{center}
\begin{tabular}{c|ccccc}
\toprule
$q$         & $0$ & $1$ & $2$ & $3$ & $4$ \\
\midrule
$\lambda$   & $(5,2)$ & $(5,2,1)$ & $(5,2,1^2)$ & $(5,2,1^3)$ & $(5,2,1^4)$ \\
$n$         & $7$     & $8$       & $9$         & $10$        & $11$ (predicted) \\
$f^{(4, 1)}$ & $4$    & $4$       & $4$         & $4$         & $4$ \\
$\dim B$    & $5$     & $8$       & $12$        & $4$         & $4$ (predicted stable) \\
factor      & $1.25$  & $2$       & $3$         & $1$         & $1$ \\
\bottomrule
\end{tabular}
\end{center}

The non-stable values $5, 8, 12$ at $q = 0, 1, 2$ admit no closed form
as multiples of $f^{(4,1)}$ alone (the factors $1.25, 2, 3$ are not all
integers), but the \emph{differences} $8 - 5 = 3$, $12 - 8 = 4$
suggest the value may be a sum of multiple $f^\mu$ contributions
across non-stable $\mu_i$'s.  We were unable to identify the precise
shape data, and leave the closed form open.

The stability threshold is sharply $q_0((5, 2)) = 3$: the dim drops
abruptly from $12$ at $q = 2$ to $4$ at $q = 3$.

\subsection{The $(k, 2)$ row at $q = 0$}

\begin{center}
\begin{tabular}{c|cccc}
\toprule
$k$        & $3$ & $4$ & $5$ & $6$ \\
\midrule
$\lambda$  & $(3, 2)$ & $(4, 2)$ & $(5, 2)$ & $(6, 2)$ \\
$n$        & $5$       & $6$       & $7$       & $8$       \\
$f^{(k-1, 1)}$ & $2$   & $3$       & $4$       & $5$       \\
$\dim B$   & --       & $3$       & $5$       & $6$       \\
factor     & --       & $1$ (stable) & $1.25$   & $1.2$     \\
\bottomrule
\end{tabular}
\end{center}

For $(k, 2)$ with $k \ge 5$, even $q = 0$ is non-stable, contrary to
the (verbal) ``the closed form starts at $q = 0$'' baseline for the
shapes $(3, 2)$ and $(4, 2)$.  This shows $q_0((k, 2))$ depends on
$k$ in a non-trivial way.

\subsection{The $(k, k, 1)$ family at $q = 1$ (row-repeat at top)}

\begin{center}
\begin{tabular}{c|cccc}
\toprule
$k$         & $2$ & $3$ & $4$ & $5$ \\
\midrule
$\lambda$   & $(2,2,1)$ & $(3,3,1)$ & $(4,4,1)$ & $(5,5,1)$ \\
$n$         & $5$       & $7$       & $9$       & $11$ \\
$f^{(k-1, k-1)}$ & $1$  & $2$       & $5$       & $14$ \\
$\dim B$    & $1$       & $3$       & $10$      & $30$ \\
factor      & $1$       & $1.5$     & $2$       & $2.14$ \\
\bottomrule
\end{tabular}
\end{center}

The factor here is $\dim V_{(k, k-1)}$ divided by $f^{(k-1, k-1)}$ at
$k = 2, 3, 4, 5$ (i.e., $1, 2, 5, 14$) divided by themselves --- so
the factor varies non-integrally.  We omit a closed form.

\subsection{The $(4, 3, 1^q)$ family across $q$}

\begin{center}
\begin{tabular}{c|ccccc}
\toprule
$q$         & $0$ & $1$ & $2$ & $3$ & $4$ \\
\midrule
$\lambda$   & $(4,3)$ & $(4,3,1)$ & $(4,3,1,1)$ & $(4,3,1^3)$ & $(4,3,1^4)$ \\
$n$         & $7$     & $8$       & $9$         & $10$        & $11$ \\
$f^{(3,2)}$ & $5$     & $5$       & $5$         & $5$         & $5$ \\
$\dim B$    & $4$     & $7$       & $12$        & $5$         & $5$ \\
factor      & $0.8$   & $1.4$     & $2.4$       & $1$         & $1$ \\
\bottomrule
\end{tabular}
\end{center}

This family illustrates the transition.  Sub-stable at $q = 0$
($\dim < f$), super-stable for $q \in \{1, 2\}$, then sharply stable
at $q \ge 3$.  The jump from $12$ at $q = 2$ to $5$ at $q = 3$ is a
\emph{phase transition}, not a gradual one.

\subsection{The $(4, 4, 1^q)$ family across $q$}

\begin{center}
\begin{tabular}{c|ccc}
\toprule
$q$         & $2$ & $3$ & $4$ (predicted) \\
\midrule
$\lambda$   & $(4,4,1,1)$ & $(4,4,1,1,1)$ & $(4,4,1^4)$ \\
$n$         & $10$        & $11$          & $12$ (compute heavy) \\
$f^{(3,3)}$ & $5$         & $5$           & $5$ \\
$\dim B$    & $10$        & $15$          & ? (likely stable, $5$, or $\ge 10$) \\
factor      & $2$         & $3$           & ? \\
\bottomrule
\end{tabular}
\end{center}

The factor \emph{grows} from $2$ to $3$ in this family.  Whether
$q = 4$ jumps back to stable (analogous to $(4, 3, 1^q)$) or continues
the growth (factor $4$?) is the leading open empirical question.
Compute infrastructure for $n = 12$ is in place but slow ($\dim
V_{(4, 4, 1^4)} = 1485$, hence $\sim 1500^2$ matrices on $11$
generators).

\subsection{Comprehensive table of measured $\dim B$ for non-stable
$\lambda$}

\begin{center}
\begin{tabular}{lccccc}
\toprule
$\lambda$ & $n$ & $r$ & $q$ & $f^{\lambda^{(\ge 2)}}$ & $\dim B$ \\
\midrule
$(3, 3)$        & $6$  & $0$ & $0$ & $2$  & $1$  (factor $0.5$) \\
$(4, 2)$        & $6$  & $1$ & $0$ & $3$  & $3$  (stable!) \\
$(5, 2)$        & $7$  & $1$ & $0$ & $4$  & $5$  (factor $1.25$) \\
$(6, 2)$        & $8$  & $1$ & $0$ & $5$  & $6$  (factor $1.2$) \\
$(4, 3)$        & $7$  & $1$ & $0$ & $5$  & $4$  (factor $0.8$) \\
$(5, 5)$        & $10$ & $0$ & $0$ & $14$ & $6$  (factor $0.43$) \\
$(3, 3, 1)$     & $7$  & $1$ & $1$ & $2$  & $3$  \\
$(4, 2, 1)$     & $7$  & $2$ & $1$ & $3$  & $6$  \\
$(3, 3, 2)$     & $8$  & $1$ & $0$ & $5$  & $3$  \\
$(4, 3, 1)$     & $8$  & $2$ & $1$ & $5$  & $7$  \\
$(5, 2, 1)$     & $8$  & $2$ & $1$ & $4$  & $8$  \\
$(3, 3, 2, 1)$  & $9$  & $2$ & $1$ & $5$  & $8$  \\
$(4, 2, 2, 1)$  & $9$  & $2$ & $1$ & $6$  & $12$ \\
$(4, 3, 1, 1)$  & $9$  & $2$ & $2$ & $5$  & $12$ \\
$(6, 2, 1)$     & $9$  & $2$ & $1$ & $5$  & $17$ \\
$(4, 4, 1, 1)$  & $10$ & $1$ & $2$ & $5$  & $10$ \\
$(4, 4, 1, 1, 1)$ & $11$ & $1$ & $3$ & $5$ & $15$ \\
$(5, 5, 1)$     & $11$ & $1$ & $1$ & $14$ & $30$ \\
$(5, 2, 1, 1)$  & $9$  & $2$ & $2$ & $4$  & $12$ \\
$(5, 2, 1^3)$   & $10$ & $2$ & $3$ & $4$  & $4$ (\emph{stable}) \\
$(6, 2, 1, 1)$  & $10$ & $2$ & $2$ & $5$  & $15$ \\
\bottomrule
\end{tabular}
\end{center}

\section{The May-13 very-late framework already explains $(4, 2, 1)$ at $q = 1$}
\label{sec:may13-framework}

The structural mechanism for non-stable $\dim B$ at \emph{one specific
family} $(4, 2, 1^q)$ was already given in the May-13 very-late paper
\texttt{2026-05-13-very-late-low-q-dim.tex} (Theorem~B, dissection at
$q = 1$).  The Phase A decomposition has \emph{four} pieces:

\begin{itemize}[topsep=0.3em,itemsep=0.2em]
\item Three corner-pair pieces $\mu_A, \mu_B, \mu_C$ (where
$\mu_C$ is the ``bad pair'' of two length-preserving corners).
\item One same-row piece $\nu_1$ (when $\lambda_1 \ge 3$).
\end{itemize}

In the \emph{stable regime}, structural hook/two-column $j$-formulas
force $\mu_C$ and $\nu_1$ to vanish (via $E_1^-$ leakage killed by a
trailing $(T_1{+}1)$ in the outer $R'$-chain), and only the good pairs
contribute.  In the \emph{non-stable regime}, the thresholds for $\mu_C$
(hook threshold $q + 1 \ge k$) or $\nu_1$ (two-column threshold) fail,
allowing the corresponding pieces to contribute non-trivially.

\paragraph{The $(4, 2, 1)$ dissection at $q = 1$.}
$\lambda = (4, 2, 1)$ at $n = 7$ has $\dim B = 6 = 2 + 1 + 1 + 2$:
\begin{itemize}[topsep=0.3em,itemsep=0.2em]
\item $\mu_A = (3, 2)$ at $n = 5$: dim $2$ (good pair).
\item $\mu_B = (4, 1)$ hook at $n = 5$: dim $1$ (good pair).
\item $\mu_C = (3, 1, 1)$ hook at $n = 5$: dim $1$ (\emph{bad pair},
hook threshold $q + 1 = 2 \ge k = 3$ fails, so $E_1^-$ leakage fails).
\item $\nu_1 = (2, 2, 1)$ same-row at row $1$, $n = 5$: dim $2$
(two-column threshold $q \ge 2$ fails at $q = 1$, so $E_1^-$ leakage
fails).
\end{itemize}

Total: $6$, matching the measured value.

\paragraph{The May-13 very-late paper does \emph{not} cover $(5, 2, 1)$ etc.}
The very-late paper proves Theorems A and B
($(3, 2, 1^q)$ at $q \ge 1$, $(4, 2, 1^q)$ at $q \ge 2$) and dissects
$(4, 2, 1)$ at $q = 1$.  It \emph{does not} establish the structural
result for $(k, 2, 1^q)$ at $k \ge 5$ in the non-stable regime, and
specifically leaves open whether the same 4-piece dissection accounts
for our new data points $(5, 2, 1)$, $(6, 2, 1)$, $(5, 2, 1, 1)$.

\section{Extending the May-13 framework to $(5, 2, 1^q)$}
\label{sec:521-extension}

We attempt the same dissection on $\lambda = (5, 2, 1)$ at $n = 8$
($\dim B = 8$).

Corners: $c_1 = (1, 5)$, $c_2 = (2, 2)$, $c_* = (3, 1)$.  Corner pairs:
\begin{align*}
   \mu_A &= (5, 2, 1) \setminus \{c_1, c_*\} = (4, 2) \text{ at } n = 6,
      & j_0(\mu_A) &= 3, \\
   \mu_B &= (5, 2, 1) \setminus \{c_2, c_*\} = (5, 1) \text{ at } n = 6,
      & j_0(\mu_B) &= 3, \\
   \mu_C &= (5, 2, 1) \setminus \{c_1, c_2\} = (4, 1, 1) \text{ at } n = 6,
      & j_0(\mu_C) &= 4.
\end{align*}

Same-row at row $1$: $\nu_1 = (3, 2, 1)$ at $n = 6$, $j_0(\nu_1) = 4$.
Rows $\ge 2$ have $\lambda_i \le 2$ so no further same-row pairs.

\paragraph{Piece contributions: computed via the commutation lemma.}
Use $B^{(\lambda)}_{j_0} V_\lambda = \mathcal{O}_\lambda \cdot \sum_X
\Phi_X(B^{(\mu_X)}_{j_0 - 1, n - 2})$ with $j_0 - 1 = 3$ on each
$\mu_X$ at $n - 2 = 6$.  We computed $\dim B^{(\mu_X)}_3$ directly
(script \texttt{test\_B\_at\_index.py}):
\begin{itemize}[topsep=0.3em,itemsep=0.2em]
\item $\mu_A = (4, 2)$ at $n = 6$: $j_0(\mu_A) = 3 = $ inner-index.
$\dim B^{(4, 2)}_3 = 3$.  Contribution upper bound $3$.
\item $\mu_B = (5, 1)$ hook at $n = 6$: $j_0(\mu_B) = 3 = $
inner-index.  Hook threshold $q + 1 \ge k$ becomes $2 \ge 5$,
\textbf{fails}.  Direct computation: $\dim B^{(5, 1)}_3 = 2$
(\emph{not} the stable hook dim $1$).  Contribution upper bound $2$.
\item $\mu_C = (4, 1, 1)$ hook at $n = 6$: $j_0(\mu_C) = 4$, inner
index $3 = j_0(\mu_C) - 1$.  Hook threshold $q + 1 \ge k$ becomes
$3 \ge 4$, \textbf{fails}.  Direct computation: $\dim B^{(4, 1, 1)}_3
= 1$ ($= \dim B^{(4, 1, 1)}_4$ also).  Contribution upper bound $1$.
\item $\nu_1 = (3, 2, 1)$ at $n = 6$: $j_0(\nu_1) = 4$, sharp.
The May-20 inclusion $B^{(3, 2, 1^{n-5})}_{\ell + 1} V \subseteq E_1^-$
is stated for $n \ge 8$, so the precise hypothesis of SRD-$421$
Theorem~$14$ is not verified at this $\nu_1$.  Empirically, we have
no immediate handle on the $\nu_1$ contribution.
\end{itemize}

\paragraph{Sum: $3 + 2 + 1 + ? = 6 + \nu_1$\textit{-contribution}.}
The corner-pair pieces alone (computed at index $3$) sum to $6$.
Measured $\dim B^{(\lambda)} = 8$.  The defect of $2$ must come from
the same-row piece $\nu_1$, or from non-additivity / loss in the
outer-chain $\mathcal{O}_\lambda$.

If we conjecture the $\nu_1$ piece contributes exactly $2$, the
4-piece dissection gives the right answer $3 + 2 + 1 + 2 = 8$.

\paragraph{Verification of the $\nu_1 = (3, 2, 1)$ inner data
(this session).}
Computing in
\texttt{\textasciitilde/projects/scratch/2026-05-14-row-repeat/test\_321\_at\_n6.py}:
\[
   \dim B^{(3, 2, 1)}_4 V_{(3, 2, 1)} \;=\; 2,
   \qquad \dim B^{(3, 2, 1)}_3 V_{(3, 2, 1)} \;=\; 2,
\]
\emph{and crucially},
\[
   \boxed{\;(T_1 + 1)\, B^{(3, 2, 1)}_4 V_{(3, 2, 1)} \ne 0
   \text{ at } n = 6\;}
\]
(the rank of $(T_1{+}1)$ applied to the $\dim$-$2$ image is also $2$,
showing $(T_1{+}1)$ is injective on it, so the image is \emph{disjoint}
from $E_1^-$).  Equivalently $B^{(3, 2, 1)}_4 V_{(3, 2, 1)} \subseteq
E_1^+$ at $n = 6$ --- the \textbf{opposite eigenspace} from the
$n \ge 8$ structural result.

This is a striking boundary phenomenon: the May-20 structural result
$B^{(3, 2, 1^{n-5})}_{\ell + 1} \subseteq E_1^-$ holds for $n \ge 8$,
but at the boundary $n = 6$ (which is $\lambda$ \emph{without} any
trailing $1$-rows beyond the necessary one), the image flips into
$E_1^+$.  Consequently SRD-$421$ Theorem~$14$'s hypothesis fails, and
the same-row $\nu_1$ piece in the $(5, 2, 1)$ decomposition is not
killed --- it contributes (conjecturally) exactly $\dim B^{(\nu_1)}_3
= 2$ to the dissection, closing the count.

\begin{conjecture}[Closed 4-piece dissection of $(5, 2, 1)$]
$\dim B^{(5, 2, 1)}_4 V_{(5, 2, 1)} = 3 + 2 + 1 + 2 = 8$, where the
four summands are the $\mu_A, \mu_B, \mu_C$ corner-pair pieces (at
their respective inner-indices in the commutation lemma) and the
$\nu_1$ same-row piece.  The $\nu_1$ piece is non-zero precisely
because $B^{(3, 2, 1)}_4 V_{(3, 2, 1)} \not\subseteq E_1^-$ at $n = 6$.
\end{conjecture}

\paragraph{Implication: the May-20 ``rank-zero'' chain has $n_0$
sharp at $n = 8$.}
The $(3, 2, 1^*)$ family was proved rank-zero in May-20 for $n \ge 8$.
This bound is \textbf{sharp}: at $n = 6$ the image is in $E_1^+$, and
in particular the rank-zero theorem (i.e.\ $\Pi^{S_n}|_{V} = 0$
through the sign-kill of $E_1^-$) cannot be derived by the same
$E_1^-$-leakage argument.  The May-13 recursive formula's claim that
``the stable closed form $f^{\lambda^{(\ge 2)}}$ matches $(3, 2, 1^q)$
at all $q \ge 1$'' is computationally correct (the dim is $2$ at
$n = 6$ too), but the structural mechanism that delivers this dim
\emph{changes} at $n = 6$: the image is in $E_1^+$ rather than
$E_1^-$.  This is a non-trivial refinement of the existing
$(3, 2, 1^q)$ structural picture.

\section{Honest assessment}

This is a failure paper: the PROVE.md conjecture is refuted, and we
have not located a correct replacement.  What we \emph{have}
established is:

\begin{itemize}[topsep=0.3em,itemsep=0.2em]
\item The two cited examples of the factor-of-two conjecture are
correct, but the conjecture as a generalisation is false at
$(6, 2, 1)$, $(3, 3, 1)$, $(4, 3, 1, 1)$, and $(4, 4, 1, 1, 1)$.
\item Non-stable $\dim B$ is a genuinely finer invariant of $\lambda$
than $f^{\lambda^{(\ge 2)}}$, with values ranging over $\{0.6, 1, 1.4,
1.5, 1.6, 2, 2.14, 2.4, 3, 3.4, \ldots\} \cdot f^{\lambda^{(\ge 2)}}$
across $15$ tested non-stable shapes.
\item The May-13 recursion fails at non-stable shapes by an
``additive defect'' that does not match the natural bad-pair
contribution; the precise correction term is open.
\end{itemize}

\paragraph{Refined open problems.}
\begin{enumerate}[topsep=0.3em,itemsep=0.2em]
\item \textbf{Pin down $q_0(\widehat\lambda)$.}  Compute
$(5, 2, 1^q)$, $(6, 2, 1^q)$, $(5, 5, 1^q)$ across $q$ to locate the
thresholds.  Is there a closed form in terms of $\widehat\lambda$?
\item \textbf{Phase transition shape.}  $(4, 3, 1^q)$ jumps from
$\dim 12$ at $q = 2$ to $\dim 5$ at $q = 3$.  Is this sharp transition
universal?  What happens in $(4, 4, 1^q)$ at $q = 4, 5, 6$?
\item \textbf{The defect formula.}  For $(4, 2, 2, 1)$ at $n = 9$,
the recursion's defect from the good-pair sum is $3$.  Is the defect
a $f^\mu$ count for some shape $\mu$?  Is there an exact
``all-pair-with-correction'' formula?
\item \textbf{Structural understanding.}  Both refined questions
above are computational; the structural question is what mechanism
permits $\dim B$ to exceed $f^{\lambda^{(\ge 2)}}$ in non-stable
regime, given that the May-15 commutation reduction would seem to
force the dim down to the stable form.  The likely answer involves
the $j$-formula leakage's $E_1^-$ assumption breaking; identifying
the exact obstruction is the natural next step.
\end{enumerate}

\paragraph{Scripts.}
All computations from
\texttt{\textasciitilde/projects/scratch/2026-05-14-row-repeat/}
(this session):
\begin{itemize}
\item \texttt{test\_base.py} --- verifies the two cited examples.
\item \texttt{scan\_small.py} --- $14$ shape scan at $n \le 10$.
\item \texttt{test\_k21.py} --- the $(k, 2, 1)$, $(k, k, 1)$, and
$(4, 3, 1^q)$ families.
\item \texttt{verify\_621.py} --- $\dim B^{(6, 2, 1)} = 17$ confirmed
at $Q \in \{11, 13, 17, 23\}$.
\end{itemize}

The underlying engine is
\texttt{\textasciitilde/projects/scratch/2026-05-13-multi-corner/compute\_dim\_B.py}.

\end{document}
